% Generated draft: requires mathematical and source review.
\documentclass[11pt]{article}
\usepackage[margin=25mm]{geometry}
\usepackage{fontspec}
\setmainfont{Latin Modern Roman}
\usepackage{amsmath,amssymb,mathtools,unicode-math}
\setmathfont{Latin Modern Math}
\usepackage{enumitem,needspace,xurl,hyperref}
\setlength{\parindent}{0pt}
\setlength{\parskip}{5pt}
\setlength{\emergencystretch}{3em}
\newcommand{\N}{\mathbb{N}}
\newcommand{\Z}{\mathbb{Z}}
\newcommand{\Q}{\mathbb{Q}}
\newcommand{\R}{\mathbb{R}}
\newcommand{\C}{\mathbb{C}}
\newcommand{\F}{\mathbb{F}}
\newcommand{\A}{\mathbb{A}}
\newcommand{\PP}{\mathbb{P}}
\newcommand{\E}{\mathbb{E}}
\newcommand{\eps}{\varepsilon}
\newcommand{\dd}{\,\mathrm{d}}
\begin{document}
\section*{UnsolvedMath ID: 20000961}
\textbf{Source problem code:} catalogue entry
\subsection{Definitions and mathematical statement}
The isolated source asks only, ``What is the minimum number of colors needed above?'' Its preceding companion record, permanent ID 20000960, ends with the assertion that every three-point collinear set $T$ can be avoided monochromatically by a $16$-coloring of every Euclidean space $\R^N$. This supplies the antecedent used by the attached aid; it is distinct from the neighboring two-color plane conjecture.

Let $T$ be any set of three distinct collinear points, with their Euclidean distances fixed. For $n\geq1$, let $\chi_n(T)$ be the least integer $q\geq1$ for which there is a coloring $c:\R^n\to\{1,\ldots,q\}$ having no monochromatic isometric image of $T$. No measurability or regularity restriction is imposed. Define the dimension- and triple-uniform number
\[
 q_* =\min\{q\geq1:\ \forall T\ \forall n\geq1\ \exists c:\R^n\to\{1,\ldots,q\}
 \text{ avoiding a monochromatic congruent }T\}.
\]
The question under this recovered antecedent is to determine $q_*$. The coloring may depend on both $T$ and $n$; it is not one coloring avoiding all triples simultaneously.

The supplied aid records $3\leq q_*\leq16$, attributing the upper bound to Straus's classical theorem and the updated lower bound to Currier--Moore--Yip's 2024 plane theorem for equally spaced triples. It also records $\chi_1(T)=2$ for each $T$ and $\chi_n(T)\leq\chi_{n+1}(T)$.

For a triple whose consecutive gaps are $ps,qs$, where $p,q$ are coprime positive integers and $s>0$, the aid gives the explicit coloring, in every finite dimension,
\[
 c(x)=\left\lfloor\frac{\|x\|_2^2}{s^2}\right\rfloor
 \pmod{(p+q)(pq+1)}.
\]
It claims this avoids every monochromatic congruent copy of that triple, yielding four colors for equal gaps and nine for ratio $1:2$. This construction applies to rational gap ratios and is not asserted for all triples. The literature bounds and coloring lemma are supplied claims requiring independent verification, not a determination of $q_*$.
\subsection{Short English statement}
Determine the least universal number of colors that can avoid any specified three-point collinear configuration in every finite Euclidean dimension. The preceding record supplies the missing antecedent. The attached aid reports bounds $3$ to $16$ and explicit dimension-uniform colorings for rational gap ratios.
\subsection{Sources}
\begin{itemize}
\item AIM workshop \emph{Graph Ramsey theory}, ``Euclidean Ramsey sets,'' question 39: \url{https://www.overleaf.com/read/mnvcscjjysvg}.
\item Local registry \texttt{data/unsolvedmath/problems.json}, permanent ID 20000961, exact source code \texttt{catalogue entry}; canonical location \texttt{aim-combinatorics-notes.json notes[85]}. Companion permanent ID 20000960, exact source code \texttt{AIM-COMBINATORICS-0085}, item 38, \texttt{notes[84]}, supplies the preceding $16$-color collinear-triple assertion.
\item The attached aid names Straus for the classical $16$-color theorem and Currier--Moore--Yip (2024) for the equally spaced triple lower bound, without full bibliographic references. Attempt 1 is labeled partial result, problem status at run open, rigor proved, and low novelty confidence; it is machine-generated, not peer reviewed, and requires independent verification before citation.
\end{itemize}
\end{document}
